\documentclass{bmstu}

\addbibresource{references.bib}

\begin{document}

\makereporttitle
{Информатика и системы управления}
{Программное обеспечение ЭВМ и информационные технологии}
{лабораторной работе №1}
{Анализ алгоритмов}
{Сравнительный анализ ресурсной эффективности алгоритмов}
{} % Вариант
{Кузнецов М. А./ИУ7-52Б} % Номер группы/ФИО студента (если несколько - разделить запятой)
{Строганов Д. В.}

\maketableofcontents

\chapter*{ВВЕДЕНИЕ}
\addcontentsline{toc}{chapter}{ВВЕДЕНИЕ}

\textbf{Цель работы} --- выполнение сравнительного анализа ресурсной эффективности алгоритмов.
\cite{placeholder-source}.

\textbf{Задачи работы}:

\renewcommand{\labelitemi}{---}
\begin{itemize}
	\item заглушка для описания решаемой задачи;
	\item заглушка для рекурсивного и нерекурсивного алгоритмов;
	\item заглушка для схем алгоритмов;
	\item заглушка для требований к программному обеспечению;
	\item заглушка для средств реализации;
	\item заглушка для реализации алгоритмов;
	\item заглушка для тестирования;
	\item заглушка для теоретической оценки памяти;
	\item заглушка для оценки временной эффективности;
	\item заглушка для сравнительной таблицы и выводов.
\end{itemize}

\chapter{Аналитическая часть}

\section{Постановка задачи}
Дан массив $A$ из $N$ чисел. Для числа $x$ требуется определить,
содержится ли $x$ в массиве $A$ и если содержится --- вернуть его
индекс в массиве.

Требуется разработать следующие алгоритмы:
\begin{itemize}
	\item линейный поиск
	\item бинарный поиск, нерекурсивная реализация, 1-й вариант
	\item бинарный поиск, нерекурсивная реализация, 2-й вариант
	\item бинарный поиск, рекурсивная реализация
\end{itemize}

Входные данные: Натуральное число $N$, массив из $N$ целых чисел,
искомое целое значение $x$. Выходные данные: индекс искомого числа в массиве
либо $-1$, если элемент отсутствует.

\chapter{Конструкторская часть}

\section{Разработка алгоритмов}
Разработаны следующие алгоритмы: линейный поиск, нерекурсивный
бинарный поиск в двух вариантах и рекурсивный бинарный поиск.

\subsubsection{Линейный алгоритм}
Нерекурсивный линейный алгоритм последовательно перебирает каждый элемент
$A_i$ массива $A$ и сравнивает его с $x$. Если текущий элемент массива равен $x$,
алгоритм возвращяет его индекс $i$ и завершается. Если ни один элемент массива 
не оказался равен $x$ --- алгоритм возвращает $-1$ и завершается.

\subsubsection{Нерекурсивный алгоритм бинарного поиска}
Нерекурсивный алгоритм итеративно ищет индекс элемента в отсортированном
массиве $A$ между позициями $l$ и $r$. Изначально позиции обозначают
края входного массива. На каждой итерации алгоритм сравнивает элемент $A_mid$ из середины диапазона с
$x$ и двигает один из краев диапазона на середину диапазона, сужая его.

Первый вариант алгоритма проверяет текущий элемент $A_mid$ c $x$ и в случае равентсва
досрочно завершает работу

\subsubsection{Рекурсивный алгоритм бинарного поиска}
Здесь будет размещена заглушка для рекурсивного варианта алгоритма.


\section{Схемы алгоритмов}
\begin{figure}[H]
	\centering
	\fbox{\rule{0pt}{45mm}\rule{0.78\textwidth}{0pt}}
	\caption{Заглушка схемы алгоритма}
	\label{fig:algorithm-placeholder}
\end{figure}

\section{Требования к программному обеспечению}
Здесь будет размещена заглушка для требований к программному обеспечению.

\chapter{Технологическая часть}

\section{Средства реализации}
Здесь будет размещена заглушка для описания языка программирования, библиотек и инструментов.

\section{Реализация алгоритмов}
Листинг реализации будет добавлен вместо этой заглушки после подготовки программного кода.

\section{Тестирование реализации алгоритмов}
\begin{table}[H]
	\centering
	\caption{Заглушка таблицы тестирования}
	\label{tab:tests-placeholder}
	\begin{tabularx}{\textwidth}{|c|X|X|X|}
		\hline
		№ & Описание теста & Входные данные & Ожидаемый результат \\ \hline
		1 & Заглушка       & Заглушка       & Заглушка            \\ \hline
	\end{tabularx}
\end{table}

Все тесты будут проверены после заполнения реализации.

\chapter{Исследовательская часть}

\section{Оценка затрачиваемой памяти}
Здесь будет размещена заглушка для оценки памяти, включая высоту дерева рекурсивных вызовов.

\section{Оценка временной эффективности}
Здесь будет размещена заглушка для оценки числа сравнений в лучшем и худшем случаях.

\subsection{Нерекурсивный алгоритм}
Здесь будет размещена заглушка для оценки временной эффективности нерекурсивного алгоритма.

\subsection{Рекурсивный алгоритм}
Здесь будет размещена заглушка для оценки временной эффективности рекурсивного алгоритма.

\section{Результаты сравнительного анализа}
\begin{table}[H]
	\centering
	\caption{Заглушка сравнительного анализа алгоритмов}
	\label{tab:comparison-placeholder}
	\begin{tabularx}{\textwidth}{|X|X|X|}
		\hline
		Критерий                & Алгоритм 1 (заглушка) & Алгоритм 2 (заглушка) \\ \hline
		Временная эффективность & Заглушка              & Заглушка              \\ \hline
		Ёмкостная эффективность & Заглушка              & Заглушка              \\ \hline
	\end{tabularx}
\end{table}

\section{Выводы}
Здесь будет размещена заглушка для выводов по критериям временной и ёмкостной эффективности.

\chapter*{ЗАКЛЮЧЕНИЕ}
\addcontentsline{toc}{chapter}{ЗАКЛЮЧЕНИЕ}

Здесь будет размещена заглушка для выводов по результатам сравнительного анализа.

\makebibliography

\appendix
\chapter{Листинг тестов}
Здесь будет размещен очевидный текст-заглушка для листингов тестов.

\end{document}
